\textbf{COLUNI 2022 -- Questão 23}\par\medskip Durante a pandemia da COVID-19, Ana realizou, em casa, um experimento para demonstrar a conversão de energia elétrica em energia térmica. Ela fez passar uma corrente elétrica, por dois minutos, através de um resistor imerso em um líquido, dentro de uma garrafa térmica. Ao final, ela verificou que a temperatura do líquido aumentou 5,2$^\circ$ C. João, um colega de Ana, relatou a ela que repetiu o seu experimento com uma corrente igual, em um resistor idêntico, pelo mesmo tempo e que observou um aumento de temperatura no líquido de 10,4$^\circ$ C. Ana formulou hipóteses para explicar a diferença nos resultados, lembrando-se de que o calor específico é uma propriedade da substância que indica a quantidade de calor que 1 g dessa substância precisa receber para aumentar a sua temperatura em 1$^\circ$ C.

\par

Das afirmativas abaixo, assinale aquela que explica CORRETAMENTE o aumento de temperatura duas vezes maior obtido por João:\begin{enumerate}
\item João utilizou a mesma massa de um líquido, com calor específico duas vezes menor.
\item João utilizou uma massa duas vezes maior de um líquido, com o mesmo calor específico.
\item João utilizou uma massa duas vezes menor de um líquido, com calor específico duas vezes maior.
\item João utilizou a mesma massa de um líquido, com calor específico duas vezes maior.\end{enumerate}
